\begin{enquestion}{20}{
    During a live-fire mortar qualification at the range, each crew fires a salvo of six rounds per mission.
    Let $X$ be the discrete random variable representing the number of direct target hits per salvo.
    Based on standard ballistic dispersion models, the theoretical his distribution is given by $P(X = x)$ in the table below.
    After monitoring $N = 250$ salvos under gusty crosswinds, the range safety officer records the following data:
    \begin{table}[htbp!]
        \centering
        \begin{tabular}{cc}
            \toprule
                \textbf{Number of hits in a salvo} & \textbf{Frequencies} & \textbf{Theoretical probability} \\
            \midrule
                $0$ & $18$ & $0.04$ \\
                $1$ & $38$ & $0.12$ \\
                $2$ & $64$ & $0.24$ \\
                $3$ & $68$ & $0.30$ \\
                $4$ & $36$ & $0.18$ \\
                $5$ & $18$ & $0.08$ \\
                $6$ & $8$  & $0.04$ \\
            \bottomrule
        \end{tabular}
    \end{table}
}

    \subquestion{3}{
        What kind of hypothesis test would be suitable for testing whether or not the wind conditions altered the hit distribution.
        Formulate the corresponding null hypothesis $H_0$ and the alternative hypothesis $H_1$ for this test.    
    }
    \solution{
        In order to test whether the wind conditions have significantly altered the hit distribution, we need to perform a chi-square goodness-of-fit test.\rubric{1}

        The null hypothesis $H_0$ and the alternative hypothesis $H_1$ can then be formulated as follows:
        \begin{itemize}
            \item[$H_0$:] the true hit distribution with the gusty crosswinds follows the theoretical ballistic model.\rubric{1}
            \item[$H_1$:] the true hit distribution with the gusty crosswinds significantly deviates from the theoretical ballistic model. \rubric{1}
        \end{itemize}
    }

    \solution{
        Om te bepalen of er een statistisch significant effect is van het al dan niet volgen van de cyber-awarenesstraining voor het adequaat reageren
        op een nep-e-mail, dienen we een chikwadraattoets voor onafhankelijkheid uit te voeren. \rubric{1}

        De nulhypothese $H_0$ en de alternatieve hypothesetoets $H_1$ kunnen dan als volgt worden geformuleerd:
        \begin{itemize}
            \item[$H_0$:] er is sprake van onafhankelijkheid tussen het wel/niet volgen van de training en het type reactie op de nep-e-mail.\rubric{1}
            \item[$H_1$:] er is GEEN sprake van onafhankelijkheid tussen het wel/niet volgen van de training en het type reactie op de nep-e-mail. \rubric{1}
        \end{itemize}
    }

    \subquestion{8}{
        Wat is de theoretische toetsingsgrootheid van deze hypothesetoets?
        Bereken de geobserveerde toetsingsgrootheid met behulp van de data in bovenstaande tabel.
    }
    \solution{
        De (theoretische) toetsingsgrootheid voor een chikwadraattoets voor onafhankelijkheid is gelijk aan
        \[
            X^2 = \sum_{i,j} \frac{(O_{ij}-E_{ij})^2}{E_{ij}}.\rubric{1}
        \]    
        Voordat we de toetsingsgrootheid kunnen berekenen, moeten we dus de verwachte frequenties $E_{ij}$ bepalen.
        Dit doen we aan de hand van de volgende formule:
        \[
            E_{ij} = \frac{(\text{rijtotaal}_i) \cdot (\text{kolomtotaal}_j)}{\text{totaal}}. \rubric{1}
        \]
        We berekenen de verwachte frequenties voor elke cel in de tabel:
        \par\medskip
        \begin{center}
            \renewcommand{\arraystretch}{1}
            \begin{tabular}{m{0.28\textwidth}|ccc|c}
                \toprule
                    & \textbf{Meldt e-mail} & \textbf{Negeert e-mail} & \textbf{Klikt op link} & \textbf{Totaal} \\
                \midrule
                \textbf{Cybertraining} & $70.1267$ & $48.1467$ & $38.7267$ & $157$ \\ 
            \textbf{Geen cybertraining} & $63.8733$ & $43.8533$ & $35.2733$ & $143$ \\ 
                \midrule
                    \textbf{Totaal} & $134$ & $92$ & $74$ & $300$ \rubric{3} \\
                \bottomrule
            \end{tabular}
        \end{center}
        \par\medskip
        Vervolgens berekenen we de geobserveerde toetsingsgrootheid $\chi^2$ als volgt:
        \begin{align*}
            \chi^2 &= \frac{(o_{1, 1} - e_{1,1})^2}{e_{1, 1}} + \frac{(o_{1, 2} - e_{1,2})^2}{e_{1, 2}} + \ldots + \frac{(o_{2, 3} - e_{2,3})^2}{e_{2, 3}} \\
                    &= \frac{(72 - 70.1267)^2}{70.1267} + \frac{(55 - 48.1467)^2}{48.1467} + \ldots + \frac{(44 - 35.2733)^2}{35.2733} \\
                    &\approx 6.2770. \rubric{3}
        \end{align*}
    }

    \subquestion{5}{
        Wat is de kansverdeling van de theoretische toetsingsgrootheid onder $H_0$? 
        Bereken met deze kansverdeling de $p$-waarde behorende bij de geobserveerde toetsingsgrootheid.
    }
    \solution{
        De theoretische toetsingsgrootheid $X^2$ volgt onder de nulhypothese een chikwadraatverdeling met $\df = (\#\text{rijen}-1)\cdot(\#\text{kolommen}-1) = (2-1)\cdot(3-1) = 2$ vrijheidsgraden.\rubric{2}

        De $p$-waarde is de kans dat de theoretische toetsingsgrootheid $X^2$ een \enquote{nog extremere} uitkomst aanneemt dan de geobserveerde toetsingsgrootheid $\chi^2$. \rubric{1}
        
        Aangezien een chikwadraattoets voor onafhankelijkheid altijd rechtszijdig is, is de $p$-waarde gelijk aan
        \begin{align*}
            p   &= P(\chi^2 > 6.2770) \\ 
                &= \chi^2\text{cdf}(\lb=6.2770; \ub=10^{99}, \df=2) \\
                &\approx 0.0433.  \rubric{2} 
        \end{align*}

        \begin{center}
            \resizebox{0.9\textwidth}{!}{
                \includegraphics{exam20261106-chi2-pvalue.png} 
            }
        \end{center}
    }

    \subquestion{4}{
        Bepaal de toetsuitslag op basis van de berekende $p$-waarde en vertaal deze terug naar de originele probleemcontext.
        Gebruik hiervoor een significantieniveau $\alpha = 0.05$.
        Geef een verklaring voor de toetsuitslag met behulp van de tabel van geobserveerde en verwachte frequenties.
    }
    \solution{
        Aangezien de $p$-waarde $p \approx 0.0433$ strikt kleiner is dan het significantieniveau $\alpha = 0.05$,
        besluiten we om de nulhypothese $H_0$ te verwerpen. \rubric{1}

        Er is voldoende reden om aan te nemen dat het gevolgd hebben van een cyber-awarenesstraining
        significant effect heeft op de manier waarop op een nep-e-mail wordt gereageerd. \rubric{1}
        
        Als we kijken naar de tabellen van geobserveerde en verwachte frequenties,
        dan zien we dat voor de cellen met kenmerk \enquote{klikt op e-mail} de groep die geen training gevolgd heeft vaker dan verwacht (onder onafhankelijkheid) op de link heeft geklikt,
        terwijl dit bij de groep die wel training gevolgd heeft andersom is.\rubric{2}
    }

    \subquestion{3}{
        State the null hypothesis $H_0$ and the alternative hypothesis $H_1$ of the hypothesis test.
        What does rejecting $H_0$ tell us about the distribution of the number of detected drones?
    }
    \solution{
        The null hypothesis $H_0$ and the alternative hypothesis $H_1$ of this test can be formulated as follows:
        \begin{align*}
            H_0&: \text{The number of detected drones follows a binomial distribution} \\
            &\qquad \qquad \text{with $n=4$ and $p=0,7$.} \rubric{1} \\
            H_1&: \text{The number of detected drones does not follow a binomial distribution} \\
            &\qquad \qquad \text{with $n=4$ and $p=0,7$.} \rubric{1}
        \end{align*}
        Rejecting the null hypothesis only tells us that the number of detected drones does not follow this specific distribution. 
        It can still be binomially distributed, but possibly with a different value for the success probability $p$. \rubric{1}
    }

    \subquestion{3}{
        Calculate the expected frequencies for each under the null hypothesis $H_0$.
    }
    \solution{
        We can evaluate the expected frequencies (under $H_0$) by using the total of $500$ independent trials and the binomial distribution with parameters $n=4$ and $p=0,7$.
        \begin{center}
            \begin{tabular}{ccc}
                \toprule
                    \textbf{Number of} & \textbf{Observed} & \textbf{Expected} \\
                    \textbf{detected drones} & \textbf{frequencies} & \textbf{frequencies}\\
                \midrule
                    $0$ & $15$  & $1000 \cdot \binompdf(n=4; p=0,7; k=0) \approx 8,1$  \\
                    $1$ & $105$ & $1000 \cdot \binompdf(n=4; p=0,7; k=1) \approx 75,6$ \\
                    $2$ & $290$ & $1000 \cdot \binompdf(n=4; p=0,7; k=2) \approx 264,6$ \\
                    $3$ & $360$ & $1000 \cdot \binompdf(n=4; p=0,7; k=3) \approx 411,6$\\
                    $4$ & $230$ & $1000 \cdot \binompdf(n=4; p=0,7; k=4) \approx 240,1$ \\
                \midrule
                    \textbf{Total:} & $1000$ & $1000$  \\ 
                \bottomrule
            \end{tabular}\rubric{3}
        \end{center}
    }

    \subquestion{9}{
        Perform the hypothesis test at a significance level $\alpha = 0,05$ by computing the $p$-value.
    }
    \solution{
        The test statistic for a chi-square goodness-of-fit test is given by:

        \begin{align*}
            X^2 &= \frac{(O_{0} - E_{0})^2}{E_{0}} + \frac{(O_{1} - E_{1})^2}{E_{1}} + \frac{(O_{2} - E_{2})^2}{E_{2}} + \frac{(O_{3} - E_{3})^2}{E_{3}} + \frac{(O_{4} - E_{4})^2}{E_{4}} \rubric{1}
        \end{align*}
        where $O_i$ and $E_i$ are respectively the observed and expected frequencies for value $i = 0, 1, 2, 3, 4$.
        We can use the table of observed and expected frequencies to calculate the observed test statistic
        \begin{align*}
            \chi^2 &= \frac{(15 - 8,1)^2}{8,1} + \ldots + \frac{(230 - 240,1)^2}{240,1} \\
                   &\approx 26,643. \rubric{3}
        \end{align*}
    
        Under the null hypothesis, the test statistic $X^2$ follows a chi-square distribution with $\text{df}= 5 - 1 = 4$ degrees of freedom. \rubric{1}
        Hence, we can calculate the $p$-value as the probability of observing a test statistic larger than $\chi^2$:
        \begin{align*}
            p = P(\chi^2 > 26,643) &= \chi^2\text{cdf}(\text{lower}=26,643; \text{upper}=10^{99}; \text{df}=4) \\
                                             &\approx 2,3471 \cdot 10^{-5}.  \rubric{2}
        \end{align*}
        As the $p$-value is (much) smaller than the significance level $\alpha$, we reject the null hypothesis.
        There is ample evidence that the number of detected drones does not follow a binomial distribution with $n=4$ and $p=0,8$. \rubric{2}
    }

    \subquestion{3}{
        Interpret the result of the hypothesis test using data from the tables of observed and expected frequencies.
    }
    \solution{
        If we consider the table of observed and expected frequencies, we see that the observed frequencies are higher than expected for $0$, $1$ or $2$ detected drones, and lower than expected for $3$ or $4$ detected drones. \rubric{2}
        It is therefore likely that the success probability of the radar is smaller than what the manufacturer claims it to be. \rubric{1}
    }

\end{enquestion}